Supervised finite-element surrogates need a solve per training example. For linear
elastostatics, minimising the assembled discrete potential energy instead is exactly
supervised regression in the stiffness norm, which measures the stress error's strain
energy, and no example is solved. Without reference solutions, the energy also ranks
predictions of a load case; its sign flags those worse than the zero field. In a
pre-registered three-dimensional comparison, a transformer trained on the energy alone beat
the same network trained on the solutions with a displacement loss on every metric, in each
of three seeds; every network supervised in that loss was worse than the zero field on some
load cases, each flagged by the energy's sign. Both replicated on held-out instances for the
stored networks. Supervised in the stiffness norm, the transformer lowered its median energy
gap by \numKdGapRel{} and von Mises error by \numKdVmRel{}, supporting two pre-registered
hypotheses, and closed its gap to the label-free network: the advantage lies in the norm of
the loss. In two dimensions, where the label-free network had the larger displacement error,
the stiffness-norm loss more than halved the supervised von Mises error, also on held-out
instances, where its energy-gap reduction did not replicate. In three dimensions, a Ritz
projection added post hoc, needing the stiffness matrix but no labels, had a smaller
displacement error than the label-free network on most instances but median peak stress
errors \numRitzValPeakRatioMin{} to \numRitzValPeakRatioMax{} times as large. No network
transferred to a substantially finer mesh without retraining.
